\begin{frame}{Pattern of Proof}
\begin{center}
\vspace{1.5cm}
{\Huge\textbf{Pattern of Proof}}
\end{center}
\end{frame}

% --------------------------------------------------
% Question 1
% --------------------------------------------------

\begin{frame}{Question 1}
A program generates an even number whenever its integer input is doubled.

The output is given by:

$$
f(n)=2n
$$

where $n$ is an integer.

Using a direct proof, show that the output of the program is always even.
\end{frame}

\begin{frame}{Question 1 -- Answer}
Let $n$ be any integer.

By the definition of an even integer, a number is even if it can be written
in the form:

$$
2k
$$

for some integer $k$.
\end{frame}

\begin{frame}{Question 1 -- Answer}
The program produces:

$$
f(n)=2n
$$

Since $n$ is an integer, $2n$ has the required form $2k$, where:

$$
k=n
$$

Therefore:

$$
\boxed{f(n)\text{ is even}}
$$

for every integer $n$.
\end{frame}

% --------------------------------------------------
% Question 2
% --------------------------------------------------

\begin{frame}{Question 2}
A data validation system accepts an integer identifier only when it is
divisible by $6$.

A developer claims that every accepted identifier must also be divisible
by $3$.

Use a direct proof to verify the developer's claim.
\end{frame}

\begin{frame}{Question 2 -- Answer}
Let $n$ be an accepted identifier.

Since the system accepts $n$ only when it is divisible by $6$, there exists
an integer $k$ such that:

$$
n=6k
$$

\end{frame}

\begin{frame}{Question 2 -- Answer}
Rewrite $6k$ as:

$$
n=3(2k)
$$

Since $k$ is an integer, $2k$ is also an integer.

Therefore, $n$ is divisible by $3$.

Hence:

$$
\boxed{6\mid n\Rightarrow3\mid n}
$$

The developer's claim is proved.
\end{frame}

% --------------------------------------------------
% Question 3
% --------------------------------------------------

\begin{frame}{Question 3}
A computer system assigns two integer values $a$ and $b$ to two processes.
Suppose the system guarantees that $a$ and $b$ have the same parity.

A monitoring program claims that $a+b$ must be even.

Prove the claim by considering the possible parity of $a$ and $b$.
\end{frame}

\begin{frame}{Question 3 -- Answer}
Since $a$ and $b$ have the same parity, there are two possible cases:

\begin{itemize}
\item both are even,
\item both are odd.
\end{itemize}
\end{frame}

\begin{frame}{Question 3 -- Answer}
\textbf{Case 1: Both are even}

Let:

$$
a=2m,\qquad b=2n
$$

Then:

$$
a+b=2m+2n
$$

$$
a+b=2(m+n)
$$

Therefore, $a+b$ is even.
\end{frame}

\begin{frame}{Question 3 -- Answer}
\textbf{Case 2: Both are odd}

Let:

$$
a=2m+1,\qquad b=2n+1
$$

Then:

$$
a+b=2m+1+2n+1
$$

$$
a+b=2(m+n+1)
$$

Therefore, $a+b$ is even.

Since both possible cases give an even result:

$$
\boxed{a+b\text{ is always even}}
$$

\end{frame}

% --------------------------------------------------
% Question 4
% --------------------------------------------------

\begin{frame}{Question 4}
A software system stores a positive integer $n$ and computes:

$$
n(n+1)
$$

A developer claims that the result is always even.

Identify an appropriate pattern of proof and use it to establish the claim.
\end{frame}

\begin{frame}{Question 4 -- Answer}
The appropriate proof pattern is a \textbf{proof by cases}.

For any integer $n$, exactly one of the following is true:

$$
n=2k
$$

or

$$
n=2k+1
$$

\end{frame}

\begin{frame}{Question 4 -- Answer}
\textbf{Case 1: $n=2k$}

Then:

$$
n(n+1)=2k(2k+1)
$$

Therefore:

$$
n(n+1)=2[k(2k+1)]
$$

So the result is even.
\end{frame}

\begin{frame}{Question 4 -- Answer}
\textbf{Case 2: $n=2k+1$}

Then:

$$
n(n+1)=(2k+1)(2k+2)
$$

Therefore:

$$
n(n+1)=2[(2k+1)(k+1)]
$$

So the result is also even.

Hence:

$$
\boxed{n(n+1)\text{ is always even}}
$$

\end{frame}

% --------------------------------------------------
% Question 5
% --------------------------------------------------

\begin{frame}{Question 5}
A distributed computing system contains $n$ processing tasks. The system
requires at least one communication link between every pair of tasks.

The number of required links is:

$$
L=\frac{n(n-1)}{2}
$$

Show that the number of links is always an integer.

Use an appropriate proof pattern based on the parity of $n$.
\end{frame}

\begin{frame}{Question 5 -- Answer}
We need to show that:

$$
\frac{n(n-1)}{2}
$$

is an integer for every positive integer $n$.

We use a proof by cases based on whether $n$ is even or odd.
\end{frame}

\begin{frame}{Question 5 -- Answer}
\textbf{Case 1: $n$ is even}

Let:

$$
n=2k
$$

Then:

$$
\frac{n(n-1)}{2}
=
\frac{2k(2k-1)}{2}
$$

Therefore:

$$
L=k(2k-1)
$$

which is an integer.
\end{frame}

\begin{frame}{Question 5 -- Answer}
\textbf{Case 2: $n$ is odd}

Let:

$$
n=2k+1
$$

Then:

$$
\frac{n(n-1)}{2}
=
\frac{(2k+1)(2k)}{2}
$$

Therefore:

$$
L=k(2k+1)
$$

which is an integer.
\end{frame}

\begin{frame}{Question 5 -- Answer}
In both possible cases, the number of links is an integer.

Therefore:

$$
\boxed{\frac{n(n-1)}{2}\in\mathbb{Z}}
$$

for every positive integer $n$.

Thus, the formula always gives a valid integer number of communication
links.
\end{frame}
